\section{Introduction}
\label{sec:introduction}
Lobbying networks to undermine public trust in climate science and delay effective policy action are extensively documented by investigative journalists and civil society organisations, including dedicated platforms such as DeSmog~\cite{desmog} and InfluenceMap's LobbyMap~\cite{lobbymap}, which publish nearly 2000 profiles of people and organisations who lobby against climate action and policy.
Despite the breadth of such work, the resulting information remains siloed across heterogeneous platforms and is published primarily in long-form prose, organised around organisations rather than around the individuals, roles, and relationships that constitute lobbying networks. 
As a consequence, even basic questions, such as who funds a given think tank, which staff previously worked for fossil fuel companies, or which board members sit across multiple front groups, require tedious manual cross-referencing across sites with incompatible structures. 
Representing this domain as a knowledge graph makes such relational knowledge explicit and queryable at scale.
% Despite the breadth of such work, information remains siloed, heterogeneous
% in format, and difficult to query or analyse computationally at scale.
% This paper presents work carried out within the ClimateSense
% project~\cite{climatesense}, which builds tools to monitor and analyse climate
% misinformation across geospatial, temporal, and topical dimensions.
% Consider an investigative journalist seeking to understand the network behind
% a prominent climate sceptic think tank: who funds it, who sits on its board,
% which staff have previously worked for fossil fuel companies, and whether those
% individuals are linked to other lobbying groups.
% The information needed exists, spread across DeSmog, LobbyMap, and Wikidata, but in incompatible formats, organised around different rationales, and
% not linked, making even this single question require hours of manual cross-referencing.
% The problem is compounded by the temporal dimension: people scatter across
% organisations, funding arrangements may be indirect, and front groups are created
% and dissolved; yet this dynamic structure is difficult to query across sources in an integrated way.

In this paper, we present LACK (Lobbying Agents, Climate, and Knowledge), a knowledge graph of persons and organisations involved in lobbying against climate science and policy, constructed using a composite pipeline combining LLM, web search, and APIs, guided by a purpose-built lightweight ontology.

Specifically, we contribute:
\begin{itemize}
    \item The \textbf{LACK ontology}, a purpose-built OWL vocabulary covering
    the key relation types in the climate lobbying domain.
    \item The \textbf{LACK knowledge graph}, comprising relations between persons and organisations, extracted from more than 2k source documents across two authoritative sources, linked to Wikidata and DBpedia,
    and published as Linked Open Data under a CC BY 4.0 licence.
%    \item An \textbf{LLM-based relation extraction pipeline}, a reusable and    evaluated method for extracting typed, temporally annotated relations from investigative journalism text.
    \item A demonstration of the \textbf{utility of LACK for real-world investigative tasks} in the climate accountability sector, validated through evaluation with domain experts from investigative journalism.
%	\item An \textbf{evaluation with domain experts} from the investigative    journalism and climate accountability sector, assessing the utility and value of the resource for real-world investigative tasks.
%    \item A \textbf{sustainability and maintenance plan}, situating LACK within
%    the ongoing ClimateSense project~\cite{climatesense} and committing to
%    continued updates as source corpora expand.
\end{itemize}

%\paragraph{Paper Structure.}
The remainder of this paper is structured as follows.
Section~\ref{sec:related} situates LACK in the related literature. Section~\ref{sec:extraction} describes the knowledge extraction pipeline, and Section~\ref{sec:ontology} presents the ontology. Section~\ref{sec:construction} covers knowledge graph construction, while Section~\ref{sec:resource} details availability, licensing, and sustainability. Section~\ref{sec:usecase} illustrates case studies, and Section~\ref{sec:evaluation} reports on pipeline evaluation and expert assessment. Section~\ref{sec:conclusion} concludes with directions for future work.